\begin{helper}[Rotated tree sequence]
Let $N$ be balanced normal globally maximum data for
$(G,k,e_0,\varphi,T)$, put $p=N.height$, and fix
$j\in\operatorname{Fin}(p)$ with $1\leq j$.  Put
$f_i=\noderef{NormalCycleEdge}(N,i)$.  Suppose that $A\ne B$ satisfy
\[
 A=\begin{cases}N.beta,&j\text{ odd},\\N.alpha,&j\text{ even},\end{cases}
 \qquad
 B=\begin{cases}N.alpha,&j\text{ odd},\\N.beta,&j\text{ even},\end{cases}
\]
and that a partial coloring $\psi$ is proper, leaves exactly $f_j$
uncolored, makes $A$ missing at $y_j$ and $B$ missing at
$y_{(j+1)\bmod p}$, satisfies
$\psi(f_i)=\varphi(f_{i+1})$ for $0\leq i<j$, and fixes every identity
outside $\{f_0,\ldots,f_j\}$.

Then there is a tree sequence $T'$ based at $f_j$ which is globally
maximum for $\psi$, has $T'.length=T.length$ and
\[
 \noderef{TreeSequenceVertexSet}(T')=
 \noderef{TreeSequenceVertexSet}(T),
\]
and has the following exact entries.  For every vertex position $r<p$,
there is $i\in\operatorname{Fin}(p)$ with
\[
 i=(j+r)\bmod p,
 \qquad T'.vertex(r)=T.vertex(i),
\]
where the latter index is cast using $p\leq T.length+1$.  Every vertex
position $r\geq p$ is unchanged.  For every edge position with $r+1<p$,
there is $i\in\operatorname{Fin}(p)$ with
\[
 i=(j+r)\bmod p,
 \qquad T'.edge(r)=f_i,
\]
and for every edge position $r\geq p-1$ there is an old edge position of
the same numeric value at which $T'.edge(r)=T.edge(r)$.
\end{helper}
\begin{proof}
The closing-cycle facts in \noderef{ScheideTrunkRotationCycle} give
$p\geq3$, injectivity of the exact identities $f_i$, and their cyclic
endpoint equations.  Define
\[
 \rho(r)=(j+r)\bmod p
\]
for $0\leq r<p$, and define the entries of $T'$ by
\[
 T'.vertex(r)=
 \begin{cases}
  T.vertex(\rho(r)),&r<p,\\
  T.vertex(r),&p\leq r\leq T.length,
 \end{cases}                                                    \tag{1}
\]
and
\[
 T'.edge(r)=
 \begin{cases}
  f_{\rho(r)},&r+1<p,\\
  T.edge(r),&p-1\leq r<T.length.
 \end{cases}                                                    \tag{2}
\]
The ranges in (1)--(2) cover all positions and agree at no overlapping
case.

We check \noderef{TreeSequence}.  Addition by $j$ modulo $p$ permutes
$\operatorname{Fin}(p)$, so the first vertex block in (1) is injective and
uses exactly the old vertices with indices below $p$.  The remaining
vertices retain their old indices.  Old vertex injectivity therefore makes
the complete new vertex map injective.

The first $p-1$ new edge identities are the cycle identities except
$f_{(j+p-1)\bmod p}$, and are pairwise distinct by the cited cycle
injectivity.  The fixed tail identities are pairwise distinct by old edge
injectivity.  They are disjoint from the cycle identities: for
$0\leq i\leq p-2$, the definition \noderef{NormalCycleEdge} identifies
$f_i$ with the old tree edge at position $i$, so old injectivity separates
it from any tail position $r\geq p-1$.  If the closing identity
$f_{p-1}$ equaled such a tail edge, its closing endpoint pair would be
$\{y_{p-1},y_0\}$, whereas the old attachment field makes the same tail
identity incident with the fresh vertex $y_{r+1}$; old vertex injectivity
contradicts $r+1\geq p$.  Thus the new edge map is injective.  Formula (2)
at $r=0$ gives $T'.edge(0)=f_j$, the required first identity.

For $r+1<p$, the cyclic endpoint equation gives
\[
 G.ends(T'.edge(r))=
 \{T'.vertex(r),T'.vertex(r+1)\},                               \tag{3}
\]
so this edge attaches the next new vertex to the immediately preceding
one.  For $r\geq p-1$, old attachment gives an old index $a\leq r$ such
that
\[
 G.ends(T.edge(r))=\{T.vertex(r+1),T.vertex(a)\}.
\]
If $a<p$, take $s=(a+p-j)\bmod p$; then $s<p\leq r+1$ and (1) gives
$T'.vertex(s)=T.vertex(a)$.  If $a\geq p$, take $s=a$.  In both cases
$s\leq r$, the other endpoint is $T'.vertex(r+1)$, and (2) supplies the
required attachment.  This proves every tree-sequence field.

We next verify \noderef{TashkinovTree}.  Delete $f_j$ from the cycle and
traverse the remaining $p-1$ identities beginning at
$y_{(j+1)\bmod p}$.  The alternating old colors from
\noderef{ScheideTrunkRotationCycle}, the parity definitions of $A,B$, and
the shift and fixed equations for $\psi$ show that every new trunk edge is
colored $A$ or $B$.  For an edge position $1\leq r<p-1$, use the new
position zero, where $A$ is missing, as the earlier witness for color $A$;
use position one, where $B$ is missing, as the witness for color $B$.
Both positions are at most $r$.

Let $r\geq p-1$.  The old Tashkinov clause, obtained from the global field
of \noderef{NormalGlobalMaximumTriple} via
\noderef{GlobalMaximumTriple}, supplies a color $c$ on $T.edge(r)$ and an
old witness position $a\leq r$ where $c$ is missing.  The tail identity is
outside the shifted cycle, by the identity-disjointness argument above, so
its $\psi$-color is still $c$.  If $c=A$ or $c=B$, use respectively the
new position zero or one, which are no later than $r$ because $p\geq3$.
Otherwise $c$ is outside the two old trunk colors, since $\{A,B\}$ is
exactly that pair.  We now verify preservation of this missing witness from
the definitions, edge by edge.  Expanding \noderef{MissingColors} and
\noderef{PresentColors}, old missingness at $T.vertex(a)$ says that there is
no exact identity $g$ incident with this vertex for which
$\varphi(g)=c$.  Suppose, for a contradiction, that $c$ were present there
under $\psi$.  By \noderef{PresentColors}, some incident exact identity $g$
would satisfy $\psi(g)=c$.  If $g$ lies outside
$\{f_0,\ldots,f_j\}$, the fixed hypothesis gives
$\varphi(g)=\psi(g)=c$, contradicting old missingness.  If $g=f_j$, the
exactly-one-uncolored hypothesis \noderef{IsOneEdgeUncolored} gives
$\psi(g)$ uncolored, again a contradiction.  Finally, if $g=f_i$ for some
$i<j$, the shift hypothesis gives
$\psi(g)=\varphi(f_{i+1})$; the alternating-color clause of
\noderef{ScheideTrunkRotationCycle} puts this value in $\{A,B\}$, contrary
to $c\notin\{A,B\}$.  These cases exhaust all identities, so $c$ is not
present under $\psi$ and hence, by \noderef{MissingColors}, remains missing
at $T.vertex(a)$.  Move this vertex to its new position by the inverse index
used in the attachment argument.  That position is still at most $r$, and
supplies the required witness.  Thus $T'$ is a Tashkinov tree.

It remains to show global maximality.  The graph-only first three clauses
of the old \noderef{KTriple} are unchanged.  The hypotheses say that
$\psi$ is proper and leaves exactly $f_j$ uncolored, so they supply the
last two clauses of a new $k$-triple.  Together with the Tashkinov property
just proved, this gives the first two clauses of
\noderef{GlobalMaximumTriple}.  Formula (1) permutes exactly the first
$p$ old vertices and fixes every later one.  By
\noderef{TreeSequenceVertexSet}, the old and new displayed vertex sets are
equal, so their cardinalities agree; the old global maximum clause gives
the required cardinality.  This proves global maximality.  Formulas
(1)--(2), read with their finite-index bounds, are exactly the four entry
formulas and the length and vertex-set equalities asserted in the
statement.
\end{proof}
